
\section{Introduction}
\label{sec:intro}

In some black-box optimization problems, objective evaluations are expensive and depend on when they are performed. Examples include aerodynamic wind-tunnel testing and thermal process control, where the underlying physical configuration $\vct{x} \in \mathbb{R}^n$ may have a fixed optimum while sensor drift or temperature fluctuations change the measurement baseline during operation. Because physical trials consume time and resources, evaluation budgets are limited, and earlier evaluations cannot be repeated under identical conditions. A related setting arises in data-driven materials discovery: in the accelerated design of magnetic high-entropy alloys, Li et al.~\cite{Li2022HEA} train machine-learning surrogates of alloy properties as functions of the composition vector $\vct{x} \in \mathbb{R}^n$ and use the surrogate scores to search the composition space. When such workflows update the surrogate sequentially with newly synthesized batches, the scores supplied to the optimization loop can vary across sequential evaluation rounds.

Motivated by such scenarios, we consider the unconstrained minimization problem
\begin{equation}
  \label{eq:min_f}
  \min_{\vct{x} \in \mathbb{R}^n} f(\vct{x}),
\end{equation}
where the latent objective function $f:\mathbb{R}^n\to\mathbb{R}$ is continuous, smooth, and bounded from below by $f_{\mathrm{low}}>-\infty$, but cannot be directly evaluated. Instead, evaluations are obtained sequentially through a drifting, noisy observation channel. Queries are indexed by a discrete time $t \in \mathbb{N}$ that increments by one with each oracle evaluation. A query issued at point $\vct{x}$ and timestamp $t$ returns
\begin{equation}
  \label{eq:oracle_intro}
  F(\vct{x},t) \;=\; \varphi(f(\vct{x}),t) + \varepsilon(\vct{x},t),
\end{equation}
where $\varphi:\mathbb{R}\times\mathbb{N}\to\mathbb{R}$ represents an unknown temporal transformation and $\varepsilon(\vct{x},t)$ is an independent zero-mean observation noise with variance bounded by $\sigma_\varepsilon^2$. When $\varphi(\cdot, t)$ is strictly increasing with respect to its first argument, the spatial minimizers of $f$ in~\eqref{eq:min_f} remain invariant across time slices. However, because evaluations are performed sequentially, observations $F(\vct{x}_i, t_i)$ collected at distinct time steps reflect disparate temporal baselines. Finite differences and static interpolation models fitted to such data can confound the spatial gradient with temporal variation.

For stationary problems where $F$ depends solely on $\vct{x}$, derivative-free optimization (DFO) is well established~\cite{ConnScheinbergVicente2009,AudetHare2017,LarsonMenickellyWild2019}. Model-based methods such as NEWUOA~\cite{Powell2006NEWUOA} and BOBYQA~\cite{Powell2009BOBYQA} construct quadratic surrogates over sparse interpolation sets, commonly with $m=2n+1$ points, by minimizing the Frobenius norm of the model Hessian change~\cite{Powell2004LeastFrobenius}. In noisy settings, Py-BOBYQA~\cite{CartisRoberts2019} provides flexible, noise-aware model management and restart strategies, while the STORM framework~\cite{ChenMenickellyScheinberg2018,BlanchetCartisMenickellyScheinberg2019} establishes global convergence under probabilistically fully-linear models. In dynamic optimization, existing methods mainly target moving optimal trajectories through gradient-based prediction--correction schemes that require temporal derivatives~\cite{Fazlyab2018,SimonettoDallAnese2020}, heuristic adaptation in evolutionary algorithms~\cite{Yazdani2025}, or temporal covariance discounting in Bayesian optimization~\cite{MoralesEnciso2015}. These approaches either use derivative information or require more samples than may be available in expensive black-box applications.

Applying classical model-based DFO to time-varying channels raises two difficulties. First, because the interpolation set is accumulated over consecutive iterations, older samples reflect earlier observation baselines; fitting a stationary polynomial surrogate to these space--time observations biases the recovered gradient. Second, the observed difference $F(\vct{x}_k, t_k) - F(\vct{x}_k+\vct{s}_k, t_k+\Delta t)$ combines latent-objective change with the channel variation accumulated over $\Delta t$. A step that improves the latent objective may therefore be rejected, while one that worsens it may be accepted. Shrinking the search radius or repeatedly re-evaluating the incumbent consumes additional queries and can reduce the local signal-to-noise ratio.

We propose TOBYQA (Time-augmented Optimization BY Quadratic Approximation), a model-based DFO algorithm that represents spatial geometry and temporal variation in one saddle-point interpolation system. Augmenting the constraint block of the least-Frobenius-norm quadratic system with a time column allows TOBYQA to recover a temporal coefficient $\gamma$ jointly with the spatial gradient and model constant from a single linear solve. Under affine temporal drift, the recovered gradient is algebraically invariant to the drift rate for any noise scale and sample radius; for general $C^2$ drift, the gradient bias is bounded by $O(L_\mu \tau_{\max}^2)$. This result motivates a drift-compensated acceptance test that subtracts the estimated temporal increment from the observed reduction. We also show that, under a rotation-invariant curvature prior, Powell's quadratic interpolation kernel is the covariance kernel of the omitted quadratic terms. This gives a generalized least-squares interpretation of the ridge regularization used to accommodate noise and relax geometric poisedness requirements. The method uses an adaptive cubic regularization scheme with an $O(n)$ closed-form step and a statistical stationarity stopping rule based on the posterior gradient covariance. Under a fully-linear model, TOBYQA achieves oracle complexity $O(\varepsilon^{-2})$. The benchmark results report higher solve rates under the tested nonstationary channels and performance comparable to classical DFO methods on stationary problems.
